\documentclass[11pt,letterpaper]{article}

\usepackage[T1]{fontenc}
\usepackage[utf8]{inputenc}
\usepackage[spanish,es-tabla,es-noshorthands]{babel}
\usepackage{lmodern}
\usepackage[margin=2.15cm]{geometry}
\usepackage{setspace}
\setstretch{1.04}
\usepackage{microtype}
\usepackage{amsmath,amssymb}
\usepackage{booktabs}
\usepackage{tabularx}
\usepackage{array}
\usepackage{float}
\usepackage{enumitem}
\usepackage{xcolor}
\usepackage[most]{tcolorbox}
\usepackage{fancyhdr}
\usepackage{titlesec}
\usepackage{csquotes}
\usepackage[style=ieee,backend=biber,sorting=none,maxbibnames=6,minbibnames=1,url=false,doi=true]{biblatex}
\usepackage[hidelinks]{hyperref}

\DeclareFieldFormat{doi}{DOI: #1}
\addbibresource{propuesta_doctoral_transformaciones_series_temporales.bib}

\definecolor{accent}{HTML}{176D68}
\definecolor{ink}{HTML}{20282C}
\definecolor{soft}{HTML}{EEF5F4}

\hypersetup{
  pdftitle={Selección automática de transformaciones de entrada para el pronóstico de series temporales},
  pdfauthor={Harvey Demian Bastidas Caicedo},
  pdfsubject={Propuesta de investigación doctoral}
}

\titleformat{\section}{\large\bfseries\color{accent}}{\thesection.}{0.6em}{}
\titleformat{\subsection}{\normalsize\bfseries\color{ink}}{\thesubsection.}{0.5em}{}
\titlespacing*{\section}{0pt}{1.0em}{0.35em}
\titlespacing*{\subsection}{0pt}{0.75em}{0.2em}

\pagestyle{fancy}
\fancyhf{}
\fancyhead[L]{\small\color{gray} Propuesta de investigación doctoral}
\fancyfoot[C]{\thepage}
\renewcommand{\headrulewidth}{0.3pt}

\setlength{\parskip}{0.22em}
\setlength{\parindent}{1.1em}
\setlist[enumerate]{leftmargin=1.5em,itemsep=0.16em,topsep=0.22em}
\setlist[itemize]{leftmargin=1.45em,itemsep=0.14em,topsep=0.22em}

\newcolumntype{L}[1]{>{\raggedright\arraybackslash}p{#1}}
\newcolumntype{Y}{>{\raggedright\arraybackslash}X}

\begin{document}

\begin{center}
{\LARGE\bfseries Selección automática de transformaciones de entrada\\[0.15em]
para el pronóstico de series temporales}\\[0.75em]
{\large Propuesta de investigación doctoral}\\[0.35em]
Harvey Demian Bastidas Caicedo
\end{center}

\vspace{0.25em}
\noindent\textbf{Palabras clave:}
series temporales; pronóstico; preprocesamiento automático; aprendizaje por transferencia; ruido; abstención.

\section*{Resumen}
\addcontentsline{toc}{section}{Resumen}

La calidad de un modelo de pronóstico depende tanto del algoritmo como de la forma en que recibe los datos.
Una transformación puede atenuar errores de medición, separar escalas o exponer relaciones temporales; también puede borrar cambios útiles, introducir retrasos o aprovechar sin querer información futura.
La práctica habitual compara algunas recetas escogidas manualmente o entrega su elección a un sistema AutoML.
Sin embargo, no existe una transformación universal: la opción conveniente depende de la variable, el proceso generador, el horizonte, el nivel de ruido, el cambio de distribución y el modelo que consumirá los datos.

Esta investigación propone aprender a seleccionar, para una tarea temporal nueva, un grafo pequeño de transformaciones de entrada que respete el orden temporal.
El sistema usará diagnósticos baratos calculados solo con datos disponibles, resultados históricos de tareas anteriores y, cuando sea necesario, evaluaciones parciales del modelo.
Podrá recomendar la siguiente configuración que conviene probar o abstenerse cuando la evidencia no permita defender una elección.
La relación señal--ruido se medirá como verdad únicamente en experimentos sintéticos o semisintéticos donde la señal y la perturbación sean conocidas; en datos reales se hablará de indicadores relacionados con ruido, no de una señal limpia supuesta.

La evaluación tendrá tres niveles.
Primero, señales de referencia y procesos multivariados con perturbaciones controladas permitirán comprobar estimación por variable, eliminación de ruido y conservación de eventos.
Segundo, bancos públicos de pronóstico medirán utilidad fuera de muestra y transferencia entre tareas completas.
Tercero, las transformaciones que superen esas pruebas se incorporarán como candidatas, no como valores predeterminados, en DOIN, una infraestructura distribuida de optimización desarrollada previamente y aplicada al dominio de negociación algorítmica.
El resultado principal será la reducción del costo para alcanzar un desempeño objetivo y del arrepentimiento frente al mejor grafo disponible bajo el mismo presupuesto.
Un resultado negativo también será concluyente: si los diagnósticos no permiten anticipar qué transformación ayuda, el método deberá pedir más evidencia o conservar la entrada original.

\section{Problema y pregunta de investigación}

Las decisiones sobre limpieza, escala, cuantización, representación espectral, separación de componentes o sincronización suelen tomarse antes de entrenar el modelo.
Ese orden las hace parecer una preparación neutral, aunque cada decisión modifica el problema de aprendizaje.
Un suavizado puede reducir el error de reconstrucción y empeorar el pronóstico; una representación en frecuencia puede ser útil para una serie periódica y perjudicial ante eventos breves; una alineación puede corregir un retraso de publicación o fabricar fuga de información si mueve observaciones futuras hacia el pasado.

Los sistemas AutoML ya seleccionan algoritmos, preprocesadores e hiperparámetros y algunos reutilizan resultados de conjuntos de datos anteriores~\cite{feurer2015autosklearn,fusi2018pmf,yang2020tensoroboe}.
Auto-FP estudia de forma explícita la búsqueda de secuencias de preprocesamiento y muestra, entre otros resultados, que la búsqueda aleatoria puede ser un competidor fuerte~\cite{qi2024autofp}.
En series temporales, las propuestas automatizadas abarcan distintas partes de la tubería y también combinan selección de arquitectura, hiperparámetros y presupuestos~\cite{meisenbacher2022review,deng2022autodl,zoller2023autosktime}.
Por tanto, la novedad no puede consistir en afirmar que el preprocesamiento se automatizará.

El vacío que se estudiará es más estrecho.
Primero, las transformaciones deben tener un contrato temporal verificable: qué datos usan para ajustarse, cuándo queda disponible cada observación y cuánto retraso introducen.
Segundo, el tratamiento puede variar por variable; aplicar el mismo filtro a canales con perturbaciones y escalas distintas puede destruir información.
Tercero, la selección debe transferirse a tareas no vistas y reconocer cuando esa transferencia es perjudicial.
Cuarto, todos los costos de diagnóstico, transformación y entrenamiento deben entrar en la comparación.

\begin{tcolorbox}[colback=soft,colframe=accent,boxrule=0.6pt,arc=1.2pt,left=6pt,right=6pt,top=5pt,bottom=5pt]
\textbf{Pregunta de investigación.}
¿Bajo qué condiciones un meta-selector, entrenado con diagnósticos baratos y evaluaciones históricas, puede escoger para una tarea no vista un grafo temporalmente válido de transformaciones de entrada, reducir el costo del pronóstico fuera de muestra frente a búsquedas sin transferencia y abstenerse cuando la recomendación podría causar transferencia negativa?
\end{tcolorbox}

\section{Objetivos e hipótesis}

\subsection{Objetivo general}

Diseñar y evaluar un método de selección automática de transformaciones de entrada para series temporales que considere diferencias entre variables, validez temporal, costo de evaluación, transferencia entre tareas y abstención ante evidencia insuficiente.

\subsection{Objetivos específicos}

\begin{enumerate}
  \item Definir un espacio pequeño y tipado de grafos de transformación, junto con un banco que distinga recuperación de una señal conocida, utilidad de pronóstico y preservación de información potencialmente útil.
  \item Desarrollar un meta-selector que use características de la tarea y resultados parciales para ordenar evaluaciones bajo presupuesto, y que produzca una recomendación solo cuando su incertidumbre lo permita.
  \item Medir la transferencia a tareas públicas no vistas y, después de superar compuertas predefinidas, evaluar la incorporación de las transformaciones seleccionadas al proceso de optimización distribuida de un dominio aplicado.
\end{enumerate}

\subsection{Hipótesis falsables}

\noindent\textbf{H1. El efecto de una transformación es condicional y predecible.}
En procesos con perturbaciones heterogéneas entre variables, los diagnósticos calculados antes del entrenamiento permitirán predecir parte de la ganancia o pérdida fuera de muestra de aplicar una transformación frente a conservar la entrada original.
La asignación por variable superará al mejor tratamiento global bajo el mismo presupuesto cuando el nivel o el tipo de perturbación difiera entre canales.
H1 falla si esa ganancia no supera un estimador que solo conoce la identidad del conjunto de datos, si desaparece al dejar familias completas fuera, o si el aparente beneficio proviene de usar información futura.

\noindent\textbf{H2. La transferencia reduce el costo de búsqueda.}
En tareas públicas no vistas, el meta-selector alcanzará un desempeño objetivo con menor costo o menor arrepentimiento simple que búsqueda aleatoria, Hyperband y BOHB, usando el mismo espacio candidato y el mismo presupuesto~\cite{li2018hyperband,falkner2018bohb}.
H2 falla si el ahorro desaparece al incluir el costo de los diagnósticos, los intentos fallidos y las evaluaciones parciales, o si no supera al mejor grafo fijo elegido solo con desarrollo.

\noindent\textbf{H3. La abstención reduce la transferencia negativa.}
Frente al mismo selector obligado a recomendar, una regla calibrada de abstención reducirá la frecuencia y magnitud de recomendaciones perjudiciales bajo cambio entre tareas, manteniendo una cobertura mínima fijada antes de la prueba.
H3 falla si la reducción se obtiene absteniéndose casi siempre, si requiere reajustar el umbral con la prueba abierta o si el riesgo entre recomendaciones supera el límite predefinido.

\section{Marco teórico y estado del arte}

\subsection{Señal, perturbación y utilidad}

En el laboratorio se observará $X_{t,j}=S_{t,j}+N_{t,j}$, donde $S$ es conocido por construcción y $N$ es una perturbación controlada.
Para cada variable $j$ se podrá calcular
\begin{equation}
  \mathrm{SNR}_j=10\log_{10}\frac{\sum_t S_{t,j}^{2}}{\sum_t N_{t,j}^{2}}.
\end{equation}
Esta definición permite evaluar si un estimador recupera el orden correcto entre niveles de ruido y si un filtro reconstruye $S$.
Las señales Blocks, Bumps, HeaviSine y Doppler ofrecen referencias clásicas con discontinuidades, oscilaciones y regularidad local distinta~\cite{donoho1994ideal,donoho1995denoising}.
Se añadirán procesos multivariados de espacio de estados, oscilaciones con frecuencia variable y sistemas no lineales para que la conclusión no dependa de una sola forma de señal.

En datos reales, $S$ y $N$ no se observan por separado.
Por eso no se reportará una ``SNR verdadera'' inferida de una serie natural.
Se calculará un vector de diagnóstico que puede incluir potencia de innovación, rugosidad local, concentración espectral, estabilidad entre escalas, faltantes, antigüedad de cada fuente y dependencia entre variables.
Cada componente deberá demostrar calibración en el laboratorio y utilidad incremental en tareas no vistas.
El criterio definitivo será el pronóstico fuera de muestra, no la apariencia suave de la serie.

Sea $D_j$ una transformación temporalmente admisible de la variable $j$ y $R_j=X_j-D_j(X_j)$ su residuo.
Se compararán tres entregas al modelo: $X$; $D(X)$; y $[X,D(X),R]$ con presupuesto de parámetros controlado.
La tercera evita asumir que todo lo removido era inútil.
Si $R$ conserva capacidad de pronóstico condicional, la transformación no podrá llamarse eliminación de ruido, aunque mejore una métrica de reconstrucción.

\subsection{Selección de pipelines y metaaprendizaje}

La selección de algoritmos parte de que una configuración no domina en todas las tareas.
Auto-sklearn combina preprocesamiento, modelo e hiperparámetros y usa desempeño previo en conjuntos similares~\cite{feurer2015autosklearn}; la factorización probabilística y TensorOboe reutilizan matrices incompletas de resultados para orientar nuevas evaluaciones~\cite{fusi2018pmf,yang2020tensoroboe}.
En pronóstico, FFORMA y FFORMPP usan características de las series para combinar modelos o anticipar su desempeño~\cite{montero2020fforma,talagala2022fformpp}, y catch22 ofrece descriptores compactos de dinámica temporal~\cite{lubba2019catch22}.

Esta propuesta no seleccionará el pronosticador completo en su primera etapa.
Congelará un conjunto pequeño de modelos aguas abajo y estudiará la decisión anterior: qué transformación o combinación de transformaciones conviene evaluar.
Tampoco supondrá una tubería estrictamente serial.
El espacio será un grafo condicionado por aplicabilidad: una descomposición común--privada requiere varias variables; una corrección de retraso requiere marcas de disponibilidad; una transformación espectral requiere una ventana histórica suficiente.
Los grafos inválidos no serán candidatos penalizados: serán configuraciones no admisibles.

\subsection{Abstención y cambio de distribución}

La opción de rechazo formaliza la decisión de no responder cuando el riesgo esperado es demasiado alto~\cite{chow1970reject}.
Aquí no significa omitir un pronóstico operativo, sino no recomendar una transformación con la evidencia disponible.
El sistema puede solicitar una evaluación adicional, elegir la identidad como alternativa conservadora o declarar que no existe soporte para transferir.
El control conforme del riesgo ofrece una referencia para calibrar pérdidas monotónicas bajo supuestos explícitos~\cite{angelopoulos2024crc}; esos supuestos se comprobarán y la validez final se medirá directamente en familias no vistas.
La no estacionariedad de las series impide tratar observaciones temporales como muestras intercambiables sin justificación~\cite{kuznetsov2016online}.

\subsection{Afirmación de novedad acotada}

AutoML puede buscar un filtro si ese filtro ya está en su espacio; no identifica por sí mismo qué parte de una serie real es ruido.
Los sistemas existentes cubren selección de preprocesadores, arquitecturas y presupuestos, mientras trabajos recientes estudian incluso la selección de ejemplos ruidosos durante el entrenamiento~\cite{taga2026faf}.
La contribución propuesta es la combinación verificable de cuatro elementos: asignación de transformaciones por variable; contratos temporales que impiden usar el futuro; transferencia entre tareas sobre un grafo de operadores; y una salida de abstención con costo completo.
La revisión sistemática del primer semestre deberá confirmar esta intersección o reducir la afirmación antes del prerregistro principal.

\section{Método}

\subsection{Objeto formal}

Una tarea $\tau$ incluye una colección de series, variables objetivo, horizonte, marcas de disponibilidad, particiones temporales y métrica.
Sea $\mathcal{G}_\tau$ el conjunto finito de grafos admisibles para esa tarea.
Para $g\in\mathcal{G}_\tau$, $L_\tau(g)$ es la pérdida de pronóstico fuera de muestra con el modelo y presupuesto congelados, y $c_\tau(g)$ es el costo total de construir diagnósticos, transformar, entrenar y evaluar.
La historia $H$ contiene tareas anteriores con sus configuraciones, costos, fallos y resultados.
El meta-selector observa descriptores $m_\tau$ disponibles antes de abrir la prueba y recomienda la siguiente evaluación, la identidad o abstención.

Bajo presupuesto $B$, el arrepentimiento simple es
\begin{equation}
 r_\tau(B)=L_\tau(\hat g_B)-\min_{g\in\mathcal{G}_\tau} L_\tau(g).
\end{equation}
Como el mínimo exige conocer todo el espacio, se calculará solo en el banco donde todas las configuraciones hayan sido evaluadas fuera del proceso de selección.
La métrica operativa complementaria será el costo hasta alcanzar una pérdida dentro de $\varepsilon$ del mejor grafo conocido.

El resultado formal previsto es una regla de certificación.
Si para todos los grafos se construyen intervalos simultáneos $[\ell_g,u_g]$ que cubren $L_\tau(g)$ con probabilidad al menos $1-\delta$, una recomendación $\hat g$ será certificable cuando
\begin{equation}
  u_{\hat g}-\min_{g\in\mathcal G_\tau}\ell_g\leq\varepsilon.
\end{equation}
La tesis estudiará el costo de alcanzar esa condición y caracterizará la región donde no puede alcanzarse con la evidencia barata.
Si dos clases de tarea producen los mismos diagnósticos observables pero tienen grafos óptimos distintos separados por más de $\varepsilon$, ningún selector limitado a esos diagnósticos puede garantizar ambas decisiones; deberá adquirir más evidencia o abstenerse.

\subsection{Espacio candidato}

La primera versión no contendrá todas las ideas posibles.
Cada operador tendrá una fase de ajuste limitada al desarrollo, una transformación reproducible, una regla de disponibilidad temporal, un costo, variables aplicables y una identidad de configuración.

\begin{table}[H]
\centering
\caption{Familias iniciales de transformaciones y condición de entrada.}
\begin{tabularx}{\textwidth}{@{}L{3.0cm}L{4.0cm}Y@{}}
\toprule
\textbf{Familia} & \textbf{Candidatos iniciales} & \textbf{Condición y control} \\
\midrule
Identidad & entrada original & Control obligatorio y alternativa ante abstención. \\
Ruido por variable & media móvil hacia atrás, mediana robusta, filtro de estado hacia adelante & Nunca se ajusta en prueba; se compara señal filtrada y residuo. \\
Resolución & cuantización escalar y no uniforme & Entra solo si el banco con ruido conocido muestra una región de resolución útil. \\
Tiempo--frecuencia & espectro o wavelet calculado sobre ventana pasada & Magnitud, fase y costo se evalúan por separado; sin ventanas centradas. \\
Común--privada & componentes entrenados en desarrollo y salida $[X,C,U]$ & Solo multivariada; no se elimina el componente común por correlación. \\
Disponibilidad & antigüedad, faltantes y retraso estimado como metadatos & No se desplazan observaciones futuras hacia el pasado. \\
\bottomrule
\end{tabularx}
\end{table}

Los detectores de patrones, el entrenamiento con corrupción del autoencoder, el enrutamiento por muestra y la asignación dinámica de anchos quedan fuera del núcleo.
Podrán convertirse en estudios posteriores solo si existen al menos dos modos útiles y una mejora oracular que justifique elegir entre ellos.

\subsection{Bancos de datos}

\textbf{Nivel 1: verdad conocida.}
Se usarán señales de referencia de denoising y procesos multivariados generados con semillas publicadas.
La cuadrícula inicial incluirá SNR por variable de $\{\infty,20,10,5,0\}$ dB y perturbaciones gaussiana blanca, temporalmente correlacionada, correlacionada entre variables, heteroscedástica, impulsiva, faltantes y retrasos.
Los niveles se aplicarán de forma homogénea y heterogénea para distinguir tratamiento global de tratamiento por variable.

\textbf{Nivel 2: datos públicos.}
El Monash Time Series Forecasting Archive aportará diversidad de dominios, frecuencias, longitudes y faltantes~\cite{godahewa2021monash}; una muestra de M4 servirá como referencia de gran escala y métricas conocidas~\cite{makridakis2020m4}.
La confirmación multivariada empleará ETT, Electricity, Traffic y Weather, conjuntos ampliamente usados en pronóstico de secuencias largas~\cite{zhou2021informer}.
La partición externa dejará fuera conjuntos o familias completas, no ventanas aleatorias de la misma serie.
Una tarea se definirá antes del análisis como una fuente de datos, variable objetivo, frecuencia y horizonte evaluados mediante orígenes temporales comunes.
Antes del experimento principal se hará un censo para fijar cuántas tareas independientes soporta cada afirmación.

\textbf{Nivel 3: aplicación distribuida.}
Los datos propios de negociación no elegirán hipótesis ni umbrales.
Después de cerrar los dos niveles anteriores, los operadores aprobados se empaquetarán como configuraciones versionadas.
La infraestructura DOIN podrá optimizar su presencia y parámetros dentro del dominio de negociación, conservar todas las evaluaciones y validar los candidatos en periodos retenidos.
El meta-selector recomendará una población inicial o una lista corta; no podrá declarar un ganador ni evitar la evaluación independiente.

\subsection{Experimentos}

\textbf{E1. Calibración de ruido por variable.}
Se estimará el perfil de perturbación usando solo el pasado y se comparará con la SNR verdadera.
Se medirán error de estimación, ordenamiento entre variables y cobertura del intervalo por tipo de ruido.
Una estimación que funciona solo en ruido gaussiano quedará etiquetada para ese régimen.

\textbf{E2. Recuperación y utilidad.}
Con el mismo conjunto de pronosticadores se compararán entrada original, filtrada, y original+filtrada+residuo.
Se medirán ganancia de SNR y error de reconstrucción donde existe $S$, pero la decisión se tomará con pérdida de pronóstico, calibración probabilística, conservación de extremos y utilidad condicional del residuo.
Se incluirán filtros centrados únicamente como techo diagnóstico con retraso declarado; nunca como alternativa desplegable.

\textbf{E3. Tratamiento global frente a tratamiento por variable.}
En perturbación heterogénea, un selector local decidirá identidad o transformación para cada canal.
Se controlará el presupuesto de parámetros y se probará si el beneficio proviene de la decisión por variable o simplemente de aumentar la dimensión de entrada.

\textbf{E4. Matriz tarea--grafo.}
Se ejecutará el espacio completo en una muestra de tareas públicas con modelos sencillos congelados: estacional ingenuo, regresión lineal sobre rezagos y un modelo neuronal pequeño.
Después se añadirá un modelo secuencial confirmatorio.
Cada registro contendrá tarea, grafo, costo, estado, semilla, partición, métricas y versiones.
Los fallos seguirán siendo observaciones del costo y la factibilidad.

Antes de abrir los resultados retenidos se publicarán el inventario de tareas, la matriz de comparaciones, la pérdida primaria, el margen de relevancia práctica y el presupuesto máximo.
Un piloto separado estimará la variabilidad y permitirá calcular la potencia mediante simulación, pero no podrá elegir la dirección ni el tamaño del efecto.
Si el número de tareas independientes no permite distinguir el margen declarado, el resultado confirmatorio será inconcluso y no se reemplazarán tareas por semillas adicionales.

\textbf{E5. Selección bajo presupuesto.}
Un vecino por características y un modelo de árboles serán los primeros meta-selectores.
Solo si los superan se ensayará un modelo más complejo.
Los comparadores serán identidad, mejor grafo fijo de desarrollo, búsqueda aleatoria, Hyperband y BOHB.
Todos usarán el mismo espacio y reloj de costo.

\textbf{E6. Abstención y cambio.}
Se calibrará una curva riesgo--cobertura en tareas de desarrollo.
En familias retenidas se comparará el selector obligado con el selector que puede pedir otra evaluación, usar identidad o abstenerse.
Se informarán cobertura, magnitud de transferencia negativa y costo adicional.

\textbf{E7. Transferencia al dominio aplicado.}
Sin cambiar el periodo retenido, se compararán el contrato vigente del dominio, el mejor grafo fijo público y la lista corta recomendada.
DOIN ejecutará la optimización con presupuesto pareado.
El resultado no será una afirmación de rentabilidad general: será una prueba de si la mejora metodológica conserva su valor al cambiar de dominio y de modelo.

\subsection{Métricas y análisis estadístico}

Para pronóstico puntual se usarán MASE y una pérdida escalada definida por conjunto~\cite{hyndman2006mase}; para pronóstico probabilístico, CRPS o pinball loss según disponibilidad.
También se medirán error en extremos, costo total, fallos, tiempo, arrepentimiento, costo hasta objetivo, cobertura y riesgo selectivo.
Las comparaciones serán pareadas sobre las mismas tareas, orígenes y semillas.
La tarea o familia retenida será la unidad de generalización; una semilla no se contará como una nueva tarea.
Se usarán intervalos por bootstrap jerárquico o bloques temporales y corrección de Holm para familias de contrastes predeclaradas.
Las comparaciones de error temporal respetarán dependencia y horizonte~\cite{diebold1995accuracy,bergmeir2018cv}.

\section{Transferencia al programa aplicado}

La investigación se conecta con el sistema existente sin convertirlo en evidencia circular.
El banco público y el contrato de operadores serán independientes de un solo pronosticador y de la infraestructura distribuida.
Los modelos base podrán reemplazarse sin cambiar los datos, las transformaciones ni las particiones.
La biblioteca de transformaciones producirá artefactos reproducibles con parámetros aprendidos únicamente en desarrollo.
La integración con DOIN ocurrirá solo después de una decisión basada en los bancos sintético y público.
El optimizador recibirá identificadores y parámetros de transformaciones aprobadas, pero volverá a evaluarlas dentro del dominio aplicado y conservará la autoridad sobre la selección final.

En el plan vigente, esta línea debe insertarse antes de ampliar búsquedas costosas de representación:
contrato temporal y banco en CPU; calibración de ruido; censo público tarea--grafo; meta-selector sencillo; abstención; integración distribuida; y solo entonces confirmación aplicada.
No se reabre una campaña terminada ni se modifica una campaña activa.
Los módulos de análisis de eventos, detección especializada, enrutamiento por muestra y asignación de capacidad mantienen planes separados.

\section{Contribuciones esperadas}

\begin{enumerate}
  \item Un banco reproducible que separa recuperación de señal, utilidad de pronóstico y conservación de información residual bajo perturbaciones por variable.
  \item Un espacio de grafos de transformación con reglas de aplicabilidad y validez temporal comprobables.
  \item Un meta-selector que ordena evaluaciones bajo presupuesto y abstiene cuando no puede certificar una recomendación dentro de un margen.
  \item Evidencia de transferencia en datos públicos y una integración condicionada con optimización distribuida en un dominio aplicado.
  \item Resultados negativos útiles: regímenes donde la identidad, una búsqueda sin transferencia o la abstención son preferibles.
\end{enumerate}

\section{Riesgos, ética y límites}

\begin{table}[H]
\centering
\begin{tabularx}{\textwidth}{@{}L{4.3cm}Y@{}}
\toprule
\textbf{Riesgo} & \textbf{Respuesta metodológica} \\
\midrule
Confundir variación real con ruido & SNR verdadera solo con componentes conocidos; en datos reales se reportan indicadores y sensibilidad. \\
Fuga de futuro & Disponibilidad y ajuste declarados por operador; pruebas que comparan salida incremental con procesamiento por lotes permitido. \\
Mejor reconstrucción, peor pronóstico & La pérdida fuera de muestra manda; se conserva y prueba el residuo. \\
AutoML ya resuelve el problema & Comparación directa con Auto-FP, búsquedas sin transferencia y sistemas temporales; novedad limitada a la intersección declarada. \\
Meta-datos insuficientes & Censo previo; si no existe soporte entre tareas, se publica el banco y no se entrena un selector complejo. \\
Costo excesivo & CPU y modelos simples primero; promoción por compuertas; límite acumulado de cómputo y energía reportado. \\
Sobreajuste al dominio financiero & Datos públicos no financieros como prueba principal; negociación algorítmica solo como aplicación posterior. \\
Impacto financiero & Sin operaciones en vivo durante la tesis experimental; simulación y cuentas de demostración bajo autorización separada. \\
\bottomrule
\end{tabularx}
\end{table}

No se requieren participantes humanos ni datos personales sensibles en el núcleo experimental.
Si se incorporan conjuntos con información humana, se usarán versiones públicas bajo sus licencias y se documentarán sus restricciones.
Los resultados, incluidos fallos y configuraciones descartadas, se publicarán con sus semillas, costos y contratos.

\section{Plan de tres años y productos}

\textbf{Año 1.}
Revisión sistemática; contrato temporal; generadores y banco de perturbaciones; operadores temporalmente válidos mínimos; estudio de SNR por variable, residual y preservación de eventos.
Producto: artículo metodológico sobre cuándo la eliminación de ruido ayuda, empata o perjudica el pronóstico.

\textbf{Año 2.}
Censo de datos públicos; matriz tarea--grafo; modelos base; meta-selector y adquisición bajo presupuesto; validación dejando familias completas fuera.
Producto: artículo sobre selección automática de transformaciones y transferencia entre tareas.

\textbf{Año 3.}
Abstención calibrada bajo cambio; confirmación multivariada; integración condicionada con DOIN y dominio aplicado; consolidación de teoría, software y tesis.
Producto: artículo sobre riesgo--cobertura y transferencia negativa, más demostrador reproducible.

\textbf{Recorte defendible.}
Si el grafo amplio resulta inabordable, la tesis conserva identidad, denoising por variable y una representación temporal adicional.
Si el meta-selector no supera modelos simples, estos se mantienen como resultado final.
Si no existe soporte para transferencia, el trabajo se cierra con el banco, el mapa de regímenes y una caracterización de imposibilidad.
El enrutamiento por observación, la adaptación continua y la selección conjunta de arquitecturas quedan como trabajo futuro.

\printbibliography[title={Referencias}]

\end{document}
